%=============================================================================
% RAPPORT PFA — INSEA / 3D SMART FACTORY
% Sujet : Conception et évaluation d'un assistant RAG documentaire
% Compilation : pdflatex rapport_PFA_RAG.tex  (deux passages recommandés)
% Bibliographie intégrée : aucun fichier .bib externe requis
%=============================================================================
\documentclass[12pt,a4paper,oneside]{report}

%------------------------------------------------------------------------------
% 1. Encodage, langue et typographie (style du rapport PFE de référence)
%------------------------------------------------------------------------------
\usepackage[utf8]{inputenc}
\usepackage[T1]{fontenc}
\usepackage[numbers,sort&compress]{natbib}
\usepackage[french]{babel}
\addto\captionsfrench{\renewcommand{\tablename}{Tableau}}
\usepackage{mathptmx}
\usepackage{microtype}
\usepackage{float}
\usepackage{setspace}
\usepackage{parskip}

%------------------------------------------------------------------------------
% 2. Mise en page INSEA : 2,5 cm latéral, 3 cm haut, 2,5 cm bas
%------------------------------------------------------------------------------
\usepackage[top=3cm,bottom=2.5cm,left=2.5cm,right=2.5cm,headheight=30pt]{geometry}
\setstretch{1.5}
\setlength{\parindent}{0pt}
\setlength{\parskip}{6pt}
\setlength{\emergencystretch}{3em}

%------------------------------------------------------------------------------
% 3. Figures, tableaux, algorithmes et code
%------------------------------------------------------------------------------
\usepackage{graphicx}
\usepackage[dvipsnames,table]{xcolor}
\usepackage{booktabs}
\usepackage{multirow}
\usepackage{tabularx}
\usepackage{array}
\usepackage{longtable}
\usepackage{amsmath,amssymb}
\usepackage{tikz}
\usetikzlibrary{arrows.meta,positioning,calc,fit,shapes.geometric,backgrounds}
\usepackage{enumitem}
\usepackage{caption}
\captionsetup{font=small,labelfont=bf,labelsep=period}
\usepackage{listings}
\usepackage[hidelinks]{hyperref}
\hypersetup{pdftitle={PFA -- Assistant RAG documentaire},pdfauthor={Elkhalil DAHANI}}

%------------------------------------------------------------------------------
% 4. En-têtes et pieds de page
%------------------------------------------------------------------------------
\usepackage{fancyhdr}
\pagestyle{fancy}
\fancyhf{}
\renewcommand{\headrulewidth}{0.4pt}
\renewcommand{\footrulewidth}{0.4pt}
\fancyhead[L]{\parbox[b]{0.64\textwidth}{\small\itshape\baselineskip=12pt
Conception et évaluation d'un assistant RAG\\
pour la documentation technique}}
\fancyhead[R]{\small\itshape PFA -- INSEA 2026}
\fancyfoot[C]{\small\thepage}

%------------------------------------------------------------------------------
% 5. Titres et code
%------------------------------------------------------------------------------
\usepackage{titlesec}
\titleformat{\chapter}[display]
  {\normalfont\Large\bfseries}
  {\chaptertitlename\ \thechapter}{12pt}{\LARGE}
  [\vspace{4pt}{\hrule height 0.5pt}\vspace{8pt}]
\titlespacing*{\chapter}{0pt}{-10pt}{20pt}
\titleformat{\section}{\normalfont\large\bfseries}{\thesection}{1em}{}
\titleformat{\subsection}{\normalfont\normalsize\bfseries}{\thesubsection}{1em}{}
\titleformat{\subsubsection}{\normalfont\normalsize\itshape}{\thesubsubsection}{1em}{}

\lstdefinestyle{pythonstyle}{
  language=Python,
  basicstyle=\ttfamily\small,
  keywordstyle=\color{MidnightBlue}\bfseries,
  commentstyle=\color{OliveGreen},
  stringstyle=\color{BrickRed},
  numbers=left,
  numberstyle=\tiny\color{gray},
  frame=single,
  rulecolor=\color{black!30},
  breaklines=true,
  showstringspaces=false,
  tabsize=4
}
\lstset{style=pythonstyle}

%------------------------------------------------------------------------------
% 6. Informations à personnaliser si nécessaire
%------------------------------------------------------------------------------
\newcommand{\NomEtudiant}{M. Elkhalil DAHANI}
\newcommand{\Filiere}{Data Science}
\newcommand{\AnneeUniversitaire}{2025--2026}
\newcommand{\EncadrantINSEA}{M. Najib OURADI}
\newcommand{\EncadrantOrganisme}{[À compléter : responsable de stage, 3D Smart Factory]}
\newcommand{\JureUn}{M. Najib OURADI}
\newcommand{\JureDeux}{M. Rachid BENMANSOUR}
\newcommand{\TitreRapport}{Conception et évaluation d'un\\assistant RAG spécialisé pour la\\documentation technique}

\newcolumntype{Y}{>{\raggedright\arraybackslash}X}

%=============================================================================
\begin{document}
%=============================================================================

%------------------------------------------------------------------------------
% PAGE DE GARDE — PFA, deux jurés INSEA exclusivement
%------------------------------------------------------------------------------
\newgeometry{top=0.8cm,bottom=1.5cm,left=2cm,right=2cm}
\begin{titlepage}
\thispagestyle{empty}
\setlength{\parindent}{0pt}

\begin{tikzpicture}[remember picture,overlay]
  \fill[green!12] ([xshift=-1cm,yshift=-1cm]current page.north west) rectangle
    ([xshift=5.5cm,yshift=-8cm]current page.north west);
  \fill[brown!45!black] ([xshift=-1cm,yshift=-1cm]current page.north west) rectangle
    ([xshift=-0.4cm,yshift=-8cm]current page.north west);
  \fill[green!12] ([xshift=1cm,yshift=0.5cm]current page.south east) rectangle
    ([xshift=-6cm,yshift=5cm]current page.south east);
  \fill[brown!45!black] ([xshift=1cm,yshift=0.5cm]current page.south east) rectangle
    ([xshift=0.4cm,yshift=5cm]current page.south east);
\end{tikzpicture}

\begin{center}
{\bfseries ROYAUME DU MAROC}\\[2pt]
{\bfseries *-*-*-*-*}\\[2pt]
{\bfseries HAUT-COMMISSARIAT AU PLAN}\\[2pt]
{\bfseries *-*-*-*-*-*-*-*}\\[2pt]
{\bfseries I\textsc{nstitut} N\textsc{ational} de S\textsc{tatistique} et d'\textsc{Économie} A\textsc{ppliquée}}\\[8pt]
\rule{0.68\textwidth}{0.5pt}\\[14pt]
{\Large\bfseries INSEA}\\[3pt]
{\large\bfseries \Filiere}
\end{center}

\vspace{2.2cm}
\begin{center}
  {\Large\bfseries Projet de Fin d'Année (PFA)}\\[4pt]
  {\bfseries *****}
\end{center}

\vspace{0.45cm}
\begin{center}
\begin{tikzpicture}
  \node[draw=brown!60!black,line width=0.6pt,fill=green!8,rounded corners=8pt,
        inner xsep=10pt,inner ysep=12pt,text width=10.5cm,align=center]
  {\bfseries\Large \TitreRapport};
\end{tikzpicture}
\end{center}

\vspace{1.1cm}
\hspace*{0.14\textwidth}%
\begin{minipage}{0.84\textwidth}
\begin{tabular}{r@{\hspace{0.45em}}l}
\underline{Préparé par} : & \textit{\textbf{\NomEtudiant}} \\[7pt]
\underline{Sous la direction de} : & \textit{\textbf{\EncadrantINSEA}} \textbf{(INSEA)} \\[2pt]
                                  & \textit{\textbf{\EncadrantOrganisme}} \\
\end{tabular}
\end{minipage}

\vspace{1.1cm}
\begin{center}
\textit{\textbf{Rapport élaboré comme exigence partielle du}}\\[9pt]
{\large\bfseries Diplôme d'Ingénieur d'État}\\[7pt]
{\bfseries \Filiere}
\end{center}

\vspace{0.65cm}
\begin{center}
\begin{minipage}{0.70\textwidth}
\textit{Devant le jury composé exclusivement de :}\\[5pt]
\textbf{\small$\blacksquare$}\quad\textit{\textbf{\JureUn}} \textbf{(INSEA)}\\[3pt]
\textbf{\small$\blacksquare$}\quad\textit{\textbf{\JureDeux}} \textbf{(INSEA)}
\end{minipage}
\end{center}

\vfill
\begin{center}
{\color{gray!70}\bfseries Année universitaire \AnneeUniversitaire}
\end{center}
\end{titlepage}
\restoregeometry

\clearpage\thispagestyle{empty}\null\clearpage

%------------------------------------------------------------------------------
% PAGES LIMINAIRES
%------------------------------------------------------------------------------
\pagenumbering{roman}
\setcounter{page}{1}

\chapter*{Dédicace}
\addcontentsline{toc}{chapter}{Dédicace}
\thispagestyle{plain}
\vspace{2cm}
Je dédie ce travail à mes parents, à ma famille et à toutes les personnes qui
ont soutenu mon parcours académique. Je l'adresse également à mes enseignants,
à mes encadrants et à mes proches, dont les conseils et les encouragements ont
accompagné la réalisation de ce projet.

\newpage
\chapter*{Remerciements}
\addcontentsline{toc}{chapter}{Remerciements}
\thispagestyle{plain}
\vspace{1.5cm}
Je remercie avant tout Dieu, le Très-Haut, pour la force et la persévérance
accordées tout au long de cette expérience.

J'adresse mes remerciements sincères à \EncadrantINSEA, pour son encadrement,
ses conseils méthodologiques et son accompagnement. Je remercie également
l'équipe de \textbf{3D Smart Factory} pour l'accueil, la disponibilité et le
cadre professionnel offerts durant le stage.

Je remercie enfin les membres du jury de l'INSEA pour l'attention portée à ce
travail, ainsi que toutes les personnes ayant contribué, directement ou
indirectement, à sa réalisation.

\newpage
\chapter*{Résumé}
\addcontentsline{toc}{chapter}{Résumé}
\thispagestyle{plain}

Ce projet de fin d'année porte sur la conception et l'évaluation d'un assistant
conversationnel fondé sur la génération augmentée par récupération
(\textit{Retrieval-Augmented Generation}, RAG) pour interroger une documentation
technique consacrée à Python, Scikit-learn et LangChain. L'objectif est de
réduire le temps de recherche documentaire tout en limitant les réponses non
soutenues par des sources.

Le système met en œuvre une collecte de documentation officielle, un nettoyage
et un découpage traçable des documents, une indexation FAISS, une recherche
hybride associant BM25 et embeddings multilingues E5, puis un reranking BGE
avant la génération par Mistral-7B-Instruct. Une interface Gradio permet à
l'utilisateur de poser des questions et de consulter les sources associées.

L'évaluation finale s'appuie sur 120 questions françaises équilibrées entre les
trois domaines, accompagnées de réponses de référence et d'identifiants de
chunks. Avec Ragas 0.4.3 et un juge Mistral, les métriques observées sont de
0,888 pour la fidélité au contexte, 0,913 pour la pertinence des réponses,
0,819 pour la précision du contexte, 0,879 pour le rappel du contexte et 0,693
pour l'exactitude factuelle. L'analyse met en évidence une faiblesse plus
marquée sur LangChain et la nécessité d'améliorer la gestion des questions hors
périmètre.

\medskip
\noindent\textbf{Mots-clés :} RAG, grands modèles de langage, recherche
hybride, FAISS, BM25, embeddings multilingues, reranking, Ragas, Mistral,
documentation technique.

\newpage
\chapter*{Abstract}
\addcontentsline{toc}{chapter}{Abstract}
\thispagestyle{plain}

This final-year project focuses on the design and evaluation of a
Retrieval-Augmented Generation (RAG) conversational assistant for querying
technical documentation on Python, Scikit-learn, and LangChain. The objective
is to reduce documentation-search time while limiting answers that are not
supported by retrieved sources.

The system implements official-documentation collection, cleaning and
traceable chunking, FAISS indexing, hybrid BM25 and multilingual E5 retrieval,
BGE reranking, and answer generation with Mistral-7B-Instruct. A Gradio
interface enables users to submit questions and inspect associated sources.

The final evaluation relies on 120 French questions balanced across the three
domains, with reference answers and reference chunk identifiers. Using Ragas
0.4.3 with a Mistral judge, the observed scores are 0.888 for faithfulness,
0.913 for answer relevancy, 0.819 for context precision, 0.879 for context
recall, and 0.693 for factual correctness. The analysis identifies a more
pronounced weakness for LangChain and highlights the need for improved
out-of-scope question handling.

\medskip
\noindent\textbf{Keywords:} RAG, large language models, hybrid retrieval,
FAISS, BM25, multilingual embeddings, reranking, Ragas, Mistral, technical
documentation.

\newpage
\tableofcontents
\listoffigures
\listoftables
\newpage

\chapter*{Liste des abréviations}
\addcontentsline{toc}{chapter}{Liste des abréviations}
\begin{longtable}{@{}p{3cm}p{10.5cm}@{}}
\toprule
\textbf{Abréviation} & \textbf{Signification}\\
\midrule
\endhead
API & Application Programming Interface.\\
BM25 & Best Matching 25, fonction de classement lexicale en recherche d'information.\\
E5 & Famille de modèles d'embeddings textuels entraînés par apprentissage contrastif faiblement supervisé.\\
FAISS & Facebook AI Similarity Search, bibliothèque de recherche de vecteurs similaires.\\
GPU & Graphics Processing Unit.\\
LLM & Large Language Model.\\
PFA & Projet de Fin d'Année.\\
RAG & Retrieval-Augmented Generation.\\
Ragas & Retrieval Augmented Generation Assessment.\\
UI & User Interface.\\
\bottomrule
\end{longtable}

\clearpage
\pagenumbering{arabic}
\setcounter{page}{1}

%=============================================================================
% INTRODUCTION GÉNÉRALE
%=============================================================================
\chapter*{Introduction générale}
\addcontentsline{toc}{chapter}{Introduction générale}
\markboth{Introduction générale}{}

Les grands modèles de langage ont profondément modifié la manière dont les
utilisateurs accèdent à l'information technique. Ils peuvent reformuler une
question et produire une réponse fluide, mais ils présentent une limite
importante : une réponse grammaticalement plausible peut rester inexacte,
incomplète ou dépourvue de source vérifiable. Cette limite est particulièrement
sensible lorsqu'un développeur recherche une précision sur une API, un paramètre
ou un comportement logiciel.

La génération augmentée par récupération propose de relier la génération à une
base documentaire externe. Le principe consiste à rechercher d'abord des
passages pertinents dans le corpus, puis à fournir ces passages au modèle avant
la génération de la réponse. L'architecture RAG formalise ainsi l'association
entre une mémoire paramétrique, portée par le modèle de langage, et une mémoire
non paramétrique, portée par l'index documentaire~\cite{lewis2020rag}.

Le présent Projet de Fin d'Année a pour objectif de concevoir un assistant RAG
spécialisé dans la documentation technique de Python, Scikit-learn et LangChain.
Le système doit récupérer des passages de documentation officielle, générer une
réponse concise et fournir les sources exploitées. Au-delà de l'implémentation,
le projet vise à établir une démarche d'évaluation traçable et reproductible.

La problématique peut être formulée ainsi :

\begin{center}
\begin{minipage}{0.88\linewidth}
\itshape
Comment concevoir un assistant RAG pour la documentation technique capable de
fournir des réponses pertinentes et justifiées par des sources, tout en
évaluant séparément la qualité de la récupération et de la génération ?
\end{minipage}
\end{center}

Pour répondre à cette question, le rapport est organisé en cinq chapitres. Le
premier présente l'organisme d'accueil et le besoin. Le deuxième décrit la
démarche de conception. Le troisième expose les fondements théoriques. Le
quatrième détaille la réalisation du système. Le cinquième présente le protocole
expérimental, les résultats Ragas et leur analyse. Enfin, une discussion met en
évidence les limites et les améliorations envisageables.

%=============================================================================
\chapter{Organisme d'accueil et contexte du projet}
\label{chap:organisme}

\section{Introduction}

Ce chapitre situe le stage dans son environnement professionnel et présente le
besoin auquel répond le projet. La réalisation d'un assistant documentaire ne
constitue pas uniquement un exercice de programmation : elle répond à une
problématique d'accès, de structuration et de fiabilité de l'information
technique.

\section{Présentation de 3D Smart Factory}

Selon les informations communiquées pour le stage, \textbf{3D Smart Factory} est
une entreprise implantée à Mohammédia, dans le quartier Lotissement La Gare. Elle
accompagne des porteurs de projets et de jeunes entrepreneurs tout au long de
leur parcours, depuis les phases de recherche et d'idéation jusqu'à la production
et à la mise sur le marché~\cite{organisationSource}. Son site de Mohammédia est
également un lieu d'échange avec un écosystème d'entrepreneurs et de partenaires.

L'accompagnement repose sur un encadrement d'experts, des formations ciblées et
la mobilisation de ressources adaptées à chaque projet. Dans l'esprit de
l'Industrie 4.0, l'objectif est de transformer une idée en solution concrète.
Les interventions portent sur des projets technologiques couvrant plusieurs
maillons de la chaîne de valeur : développement, fabrication et
commercialisation~\cite{organisationSource}. Le support transmis mentionne aussi
un ancrage sectoriel dans les dispositifs médicaux et un fonctionnement de type
TPE, favorable à un accompagnement personnalisé et réactif.

\begin{figure}[H]
\centering
\begin{tikzpicture}[node distance=1.6cm,>=Latex]
\node[draw,brown!60!black,rounded corners=5pt,fill=green!8,minimum width=3cm,
      minimum height=1cm,align=center] (idea) {Idéation et\\cadrage};
\node[draw,brown!60!black,rounded corners=5pt,fill=green!8,minimum width=3cm,
      minimum height=1cm,align=center,right=of idea] (proto) {Prototypage et\\développement};
\node[draw,brown!60!black,rounded corners=5pt,fill=green!8,minimum width=3cm,
      minimum height=1cm,align=center,right=of proto] (market) {Accompagnement\\vers le marché};
\draw[->,thick] (idea) -- (proto);
\draw[->,thick] (proto) -- (market);
\node[above=0.7cm of proto,font=\bfseries] {Logique d'accompagnement de 3D Smart Factory};
\end{tikzpicture}
\caption{Synthèse du positionnement de l'organisme d'accueil, d'après les informations fournies pour le stage.}
\label{fig:3dsmartfactory}
\end{figure}

\section{Organisation générale et intérêt du cadre de stage}

L'organisation repose sur une logique d'accompagnement intégrée : une idée est
progressivement transformée en projet opérationnel grâce à un appui technique,
méthodologique et stratégique. Cette approche favorise une collaboration avec
des profils complémentaires et une confrontation rapide entre les choix
techniques et les contraintes d'usage.

Dans le cas présent, le sujet RAG s'inscrit dans une démarche de valorisation
de connaissances techniques. Un assistant documentaire peut faciliter la
consultation de ressources complexes, améliorer la productivité lors de la
recherche d'information et soutenir l'appropriation de technologies modernes
par des utilisateurs non spécialistes.

\section{Contexte, besoin et périmètre}

La documentation officielle est généralement riche, mais volumineuse et
répartie entre de nombreuses pages. Une question technique précise peut exiger
la consultation de plusieurs sections, tandis que les moteurs de recherche
classiques privilégient parfois une correspondance lexicale partielle. En
parallèle, un modèle de langage généraliste peut répondre sans disposer de la
version exacte de la documentation.

Le besoin identifié consiste donc à concevoir un système qui aide l'utilisateur
à interroger un corpus documentaire limité et identifiable. Le périmètre du
projet couvre les trois domaines suivants : Python, Scikit-learn et LangChain.
Pandas, Django, React ou d'autres bibliothèques non collectées ne font pas partie
du corpus et ne doivent pas recevoir de réponse factuelle non sourcée.

\begin{table}[H]
\centering
\caption{Exigences fonctionnelles et non fonctionnelles du projet.}
\label{tab:exigences}
\begin{tabularx}{\textwidth}{p{3.5cm}Y Y}
\toprule
\textbf{Catégorie} & \textbf{Exigence} & \textbf{Réponse apportée}\\
\midrule
Fonctionnelle & Poser une question technique. & Interface Gradio et fonction \texttt{pipeline.answer(question)}.\\
Fonctionnelle & Retrouver des preuves documentaires. & Recherche hybride et affichage des sources.\\
Fonctionnelle & Produire une réponse exploitable. & Génération Mistral guidée par le contexte extrait.\\
Qualité & Limiter les réponses hors source. & Prompt contraint, provenance des chunks et garde hors périmètre.\\
Qualité & Pouvoir reproduire l'expérience. & Notebook Kaggle, manifestes, fichiers JSON/CSV et versions du code.\\
Sécurité & Ne pas exposer les clés API. & Utilisation de Kaggle Secrets pour les juges Ragas.\\
\bottomrule
\end{tabularx}
\end{table}

\section{Objectifs du projet}

Le projet poursuit trois objectifs complémentaires. Le premier est technique :
construire un pipeline RAG complet, depuis la collecte de documents jusqu'à la
génération d'une réponse. Le deuxième est applicatif : fournir une interface
simple pour démontrer l'utilisation du système. Le troisième est scientifique :
mesurer distinctement la qualité des contextes récupérés et la qualité des
réponses générées.

\section*{Conclusion du chapitre}
\addcontentsline{toc}{section}{Conclusion du Chapitre 1}

L'environnement de 3D Smart Factory offre un cadre orienté vers l'innovation et
la mise en œuvre de solutions concrètes. Le besoin retenu est celui d'un accès
fiable à une documentation technique ciblée. Le chapitre suivant présente la
démarche adoptée pour transformer ce besoin en architecture RAG reproductible.

%=============================================================================
\chapter{Démarche scientifique et conception de l'architecture}
\label{chap:conception}

\section{Méthodologie de réalisation}

La réalisation a été structurée en étapes successives afin de séparer la qualité
des données, la qualité de la récupération, la génération et l'évaluation. Cette
séparation permet d'identifier plus facilement l'origine d'une erreur : une
réponse faible peut provenir d'un corpus incomplet, d'une récupération mal
classée ou d'une génération imprécise.

\begin{table}[H]
\centering
\caption{Étapes de la démarche de réalisation.}
\label{tab:demarche}
\begin{tabularx}{\textwidth}{p{1.2cm}p{3.2cm}Y}
\toprule
\textbf{Étape} & \textbf{Intitulé} & \textbf{Production principale}\\
\midrule
1 & Collecte & Documents issus des dépôts et documentations officielles.\\
2 & Nettoyage & Documents nettoyés, dédoublonnés et accompagnés de métadonnées.\\
3 & Benchmarking & Paramètres de recherche testés sur les questions de développement.\\
4 & Chunking et indexation & Chunks identifiés et index FAISS reproductible.\\
5 & Recherche et génération & Réponse, sources et contexte exact inséré dans le prompt.\\
6 & Évaluation & Artefacts Ragas, scores par question et manifestes d'exécution.\\
\bottomrule
\end{tabularx}
\end{table}

\section{Architecture fonctionnelle}

L'architecture est organisée autour de deux flux. Le premier est un flux hors
ligne : il prépare le corpus et construit l'index. Le second est un flux en
ligne : il traite la question de l'utilisateur, récupère des passages et génère
une réponse. Cette organisation évite de répéter la collecte ou l'indexation à
chaque question.

\begin{figure}[H]
\centering
\begin{tikzpicture}[node distance=0.8cm and 0.9cm,>=Latex,
  block/.style={draw,rounded corners=4pt,fill=green!8,minimum width=2.55cm,
  minimum height=0.8cm,align=center,font=\small},
  data/.style={draw,rounded corners=4pt,fill=blue!6,minimum width=2.55cm,
  minimum height=0.8cm,align=center,font=\small}]
\node[data] (docs) {Documentation\\officielle};
\node[block,right=of docs] (clean) {Nettoyage et\\dédoublonnage};
\node[block,right=of clean] (chunk) {Chunks avec\\provenance};
\node[data,right=of chunk] (index) {FAISS +\\BM25};
\node[block,below=1.35cm of index] (retrieve) {Recherche hybride\\top 20};
\node[block,left=of retrieve] (rerank) {Reranking BGE\\top 5};
\node[block,left=of rerank] (guard) {Contrôle de\\périmètre};
\node[block,left=of guard] (question) {Question\\utilisateur};
\node[block,below=1.35cm of rerank] (llm) {Mistral-7B\\génération};
\node[data,right=of llm] (answer) {Réponse +\\sources};
\draw[->,thick] (docs)--(clean);
\draw[->,thick] (clean)--(chunk);
\draw[->,thick] (chunk)--(index);
\draw[->,thick] (index)--(retrieve);
\draw[->,thick] (retrieve)--(rerank);
\draw[->,thick] (rerank)--(guard);
\draw[->,thick] (question)--(guard);
\draw[->,thick] (guard)--(llm);
\draw[->,thick] (llm)--(answer);
\draw[->,dashed] (guard.south) to[bend right=25] node[left,font=\scriptsize,align=center]{hors\\périmètre} (answer.south);
\end{tikzpicture}
\caption{Architecture fonctionnelle de l'assistant RAG.}
\label{fig:architecture}
\end{figure}

\section{Choix des données et traçabilité}

Le corpus est constitué à partir de documentations techniques officielles. Ce
choix réduit la dépendance à des contenus non vérifiés et facilite la citation
des sources. Chaque chunk possède un identifiant stable, l'identifiant du
document, l'URL d'origine, la version de la source et une empreinte de contenu.
Cette information permet de relier une réponse, un passage utilisé et une ligne
du jeu d'évaluation.

Le dédoublonnage est appliqué physiquement aux documents nettoyés, afin que le
nombre de documents annoncé corresponde réellement au corpus indexé. Cette
précaution est importante pour la reproductibilité : une différence entre le
corpus déclaré et le corpus interrogé rendrait les résultats difficiles à
interpréter.

\section{Stratégie de récupération}

La version finale associe une recherche dense et une recherche lexicale. Le
modèle \texttt{intfloat/\allowbreak multilingual-e5-\allowbreak base} encode la question française avec
le préfixe \texttt{query:} et les passages avec le préfixe \texttt{passage:}.
Ce choix est adapté à la présence de questions françaises et de documentation
principalement anglaise. Les embeddings de phrases permettent de représenter et
comparer des contenus sémantiquement similaires~\cite{reimers2019sbert,wange5}.

La recherche BM25 conserve une capacité utile pour les noms d'API et les termes
techniques exacts. Les scores dense et lexical sont combinés avant de sélectionner
vingt candidats. Un reranker cross-encoder BGE réordonne ensuite ces candidats
et conserve les cinq passages remis au générateur. La séparation entre candidats
et contextes finaux limite le bruit transmis au modèle de génération.

\section{Critères de conception de l'évaluation}

Une baseline historique de 20 questions est conservée comme diagnostic initial.
L'évaluation principale s'appuie cependant sur 120 questions : 40 pour Python,
40 pour Scikit-learn et 40 pour LangChain. Chaque enregistrement comprend une
question française, une réponse de référence, des URLs documentaires et un ou
plusieurs identifiants de chunks de référence. Ce dispositif rend possible une
analyse de récupération indépendante du jugement automatique.

\section*{Conclusion du chapitre}
\addcontentsline{toc}{section}{Conclusion du Chapitre 2}

La démarche retenue sépare clairement préparation des données, récupération,
génération et évaluation. Cette séparation constitue une condition essentielle
pour diagnostiquer le système et améliorer une composante sans masquer les
limites d'une autre. Le chapitre suivant précise les fondements théoriques qui
justifient ces choix.

%=============================================================================
\chapter{Fondements théoriques du RAG et de son évaluation}
\label{chap:theorie}

\section{Objectif et organisation du chapitre}

Ce chapitre présente les notions nécessaires pour comprendre l'architecture et
les résultats du projet. Il distingue le rôle du modèle de langage, la recherche
d'information, la sélection des passages et l'évaluation. Cette séparation est
essentielle : dans un système RAG, une réponse faible peut provenir du corpus,
du récupérateur, du classement ou de la génération. Les revues récentes du
domaine confirment que l'évaluation doit couvrir ces composantes, plutôt que de
noter uniquement la réponse finale~\cite{gao2023survey,yu2024eval}.

\section{Grands modèles de langage et génération conditionnée}

Un grand modèle de langage prédit la suite la plus probable d'une séquence de
tokens en s'appuyant sur des régularités apprises durant son entraînement. Pour
une sortie $y=(y_1,\ldots,y_T)$ et une entrée $x$, cette génération peut être
représentée par :
\[
P_\theta(y\mid x)=\prod_{t=1}^{T}P_\theta(y_t\mid y_{<t},x),
\]
où $\theta$ représente les paramètres appris. Les architectures Transformer
fondées sur l'attention constituent la base de ces modèles~\cite{vaswani2017attention}.
Les encodeurs bidirectionnels de type BERT sont, eux, particulièrement adaptés à
la représentation de phrases et au classement de paires de textes~\cite{devlin2018bert}.

Cette capacité de génération ne garantit pas une connaissance exacte, récente ou
vérifiable d'une API. Un LLM peut produire une réponse cohérente en surface sans
qu'elle soit soutenue par une documentation. Le projet utilise
\texttt{Mistral-7B-Instruct-v0.3} comme générateur local, chargé par la
bibliothèque Transformers~\cite{jiang2023mistral,wolf2019transformers}. La
quantification 4-bit est un choix d'ingénierie permettant son exécution sur le
GPU T4 de Kaggle ; elle ne transforme pas le modèle en source documentaire.

\section{Principe de la génération augmentée par récupération}

Le RAG combine une mémoire paramétrique, intégrée dans le modèle de langage, et
une mémoire externe consultable au moment de la question. Lewis et al.
montrent que cette association permet d'utiliser un index de passages pour
améliorer la génération sur des tâches nécessitant des connaissances
spécifiques~\cite{lewis2020rag}.

Dans le cas d'une documentation technique, la mémoire externe présente deux
avantages majeurs. Elle peut être mise à jour sans réentraîner le modèle et
elle fournit une provenance pour les réponses. Néanmoins, l'amélioration de la
génération dépend directement de la qualité des passages récupérés : un contexte
hors sujet peut conduire à une réponse hors sujet.

\section{Préparation du corpus : nettoyage, chunks et provenance}

Une documentation complète ne peut pas être transmise intégralement au modèle
pour chaque question. Elle est donc découpée en fragments, appelés \textit{chunks}.
Un chunk doit être assez court pour être traité efficacement, tout en conservant
un contexte sémantique utile. Un découpage trop court sépare une définition de
ses conditions d'usage ; un découpage trop long dilue la précision de la
recherche et augmente le bruit dans le prompt.

Le projet conserve, pour chaque chunk, un identifiant stable, l'identifiant du
document d'origine, l'URL, le domaine et une empreinte de contenu. Ces
métadonnées permettent d'associer une réponse à une source et de vérifier, dans
le jeu de test, si le passage annoté a effectivement été retrouvé. Le nettoyage
et le dédoublonnage sont donc des opérations méthodologiques : ils améliorent la
cohérence de l'index et la reproductibilité du protocole.

\section{Embeddings, similarité vectorielle et FAISS}

Un embedding transforme un texte en vecteur numérique de dimension fixe. La
proximité entre une question $q$ et un passage $c$ peut être mesurée par la
similarité cosinus :
\[
\operatorname{sim}(q,c)=\frac{\mathbf{v}_q\cdot\mathbf{v}_c}
{\lVert\mathbf{v}_q\rVert\,\lVert\mathbf{v}_c\rVert}.
\]
Sentence-BERT a popularisé cette approche en produisant des embeddings de
phrases efficaces pour la recherche sémantique~\cite{reimers2019sbert}. Les
récupérateurs denses bi-encodeurs ont également démontré qu'un encodage séparé
des questions et passages est utilisable pour les questions-réponses ouvertes~\cite{karpukhin2020dpr}.

Le modèle E5 retenu dans le projet est entraîné par apprentissage contrastif
faiblement supervisé et vise des représentations génériques de textes~\cite{wange5}.
Les préfixes \texttt{query:} et \texttt{passage:} sont appliqués conformément à
son protocole. Ils conviennent au cas étudié : les questions sont majoritairement
en français alors que les documentations techniques sont principalement anglaises.

FAISS est utilisé pour indexer les vecteurs des chunks et retrouver rapidement
les voisins les plus proches d'une question encodée. Les travaux de Johnson,
Douze et Jégou décrivent des techniques de recherche de similarité à grande
échelle adaptées aux vecteurs de forte dimension~\cite{johnson2017faiss}. Dans
le présent projet, FAISS assure le premier niveau de récupération dense.

\section{Recherche BM25 et recherche hybride}

BM25 est une fonction de classement issue de la recherche d'information. Elle
valorise notamment la présence de termes de la requête dans un document, tout en
corrigeant la fréquence des mots et la longueur des documents~\cite{robertson2009bm25}.
Cette approche reste particulièrement pertinente lorsqu'une question contient
un nom de fonction, une classe ou un paramètre exact.

Pour une requête $q$ et un document $d$, la fonction BM25 peut s'écrire :
\[
\operatorname{BM25}(q,d)=\sum_{t\in q}\operatorname{IDF}(t)
\frac{f(t,d)(k_1+1)}{f(t,d)+k_1\left(1-b+b\frac{|d|}{\operatorname{avgdl}}\right)},
\]
où $f(t,d)$ est la fréquence du terme $t$ dans $d$, $|d|$ la longueur du
document et $\operatorname{avgdl}$ la longueur moyenne des documents.

La recherche dense capte davantage le sens général, tandis que BM25 favorise
les correspondances lexicales. Leur combinaison est appelée recherche hybride.
Elle est utile pour des questions techniques, car celles-ci mélangent souvent
un concept exprimé en langage naturel et un symbole ou une API précise.

Dans le système final, les scores normalisés sont fusionnés selon :
\[
S_{\mathrm{hybride}}(q,c)=\alpha S_{\mathrm{dense}}(q,c)+(1-\alpha)S_{\mathrm{BM25}}(q,c),
\quad \alpha=0{,}7.
\]
Ce poids donne une place majoritaire à la proximité sémantique tout en conservant
le bénéfice des termes exacts. Il s'agit d'un paramètre choisi lors du
benchmarking du projet et non d'une valeur universelle applicable à tout corpus.

\section{Reranking par cross-encoder}

La recherche initiale privilégie la vitesse : elle doit examiner un grand nombre
de chunks. Le reranking intervient ensuite sur un ensemble réduit de candidats.
Un cross-encoder évalue conjointement la question et un passage et peut donc
produire un classement plus fin qu'une simple comparaison de vecteurs. Cette
approche à deux étapes est courante pour améliorer la précision finale de la
recherche~\cite{nogueira2019passagererank}.

Le projet utilise \texttt{BAAI/bge-reranker-v2-m3}, un reranker multilingue~\cite{bgereranker}. Il
réordonne les vingt candidats hybrides et transmet les cinq meilleurs passages
au générateur. Ce paramètre \textit{top 5} est également réutilisé dans
l'évaluation des contextes.

\section{Hallucination et contrôle du périmètre}

Une hallucination désigne ici une information non soutenue par les documents
fournis ou une réponse donnée alors que le corpus ne couvre pas le sujet. Le
prompt impose donc au générateur de s'appuyer uniquement sur les passages
extraits et d'indiquer une insuffisance documentaire lorsque les passages ne
permettent pas de répondre.

Cette instruction reste utile mais insuffisante. Le projet ajoute en conséquence
un contrôle déterministe de périmètre. Si une bibliothèque explicitement absente
du corpus est identifiée, comme \texttt{pandas} ou l'alias \texttt{pd.}, le
pipeline renvoie un refus avant la récupération et avant l'appel au modèle
Mistral. Cette décision réduit à la fois le risque de réponse inventée et le
coût de calcul.

\section{Évaluation Ragas}

Ragas propose des métriques destinées à mesurer plusieurs dimensions d'un
système RAG, notamment la récupération et la fidélité de la réponse au
contexte~\cite{es2023ragas,ragasdocs}. Dans le protocole final, les métriques sont
calculées avec un juge Mistral externe. Elles sont complétées par des mesures
déterministes fondées sur les identifiants de chunks de référence. Les travaux de
synthèse sur l'évaluation RAG confirment l'importance de distinguer la qualité
de la récupération de la qualité de la génération~\cite{yu2024eval}.

\begin{table}[H]
\centering
\caption{Interprétation des métriques utilisées dans l'évaluation finale.}
\label{tab:metriques}
\begin{tabularx}{\textwidth}{p{3.4cm}Y}
\toprule
\textbf{Métrique} & \textbf{Interprétation dans le projet}\\
\midrule
Faithfulness & Vérifie si les affirmations de la réponse sont soutenues par les contextes remis au générateur.\\
Answer Relevancy & Vérifie si la réponse traite effectivement la question posée.\\
Context Precision & Évalue si les passages récupérés sont utiles pour répondre.\\
Context Recall & Évalue si les passages récupérés couvrent les informations nécessaires.\\
Factual Correctness & Compare la réponse générée à la réponse de référence.\\
Hit@5 et MRR@5 & Mesures déterministes indiquant la présence et le rang du chunk de référence dans les cinq passages récupérés.\\
\bottomrule
\end{tabularx}
\end{table}

Les métriques Ragas du projet sont calculées par un juge LLM. Elles doivent donc
être lues comme des évaluations automatiques structurées, mais non comme une
vérité absolue. Elles sont complétées par deux mesures déterministes :
\[
\operatorname{Hit@}k=\frac{1}{N}\sum_{i=1}^{N}\mathbb{1}[R_i\cap C_i^{(k)}\neq\varnothing],
\qquad
\operatorname{MRR@}k=\frac{1}{N}\sum_{i=1}^{N}\frac{1}{\operatorname{rang}_i}.
\]
Ici, $R_i$ représente les chunks de référence de la question $i$, $C_i^{(k)}$
les $k$ passages récupérés et $\operatorname{rang}_i$ le rang du premier chunk
de référence, lorsqu'il est présent. Ces indicateurs vérifient directement la
présence et la position des preuves annotées.

\section*{Conclusion du chapitre}
\addcontentsline{toc}{section}{Conclusion du Chapitre 3}

Les fondements théoriques montrent que la qualité d'un système RAG dépend d'une
chaîne complète : données de qualité, découpage traçable, récupération dense et
lexicale, classement, contexte, génération contrôlée et évaluation à plusieurs
niveaux. Le chapitre suivant traduit ces principes en choix d'implémentation dans
le code et le notebook Kaggle.

%=============================================================================
\chapter{Réalisation du système RAG}
\label{chap:realisation}

\section{Environnement de développement}

Le système est développé en Python et exécuté principalement dans Kaggle. Cet
environnement fournit un GPU T4 compatible avec la quantification 4-bit du
modèle génératif. Le code est organisé en scripts correspondant aux étapes du
pipeline, complétés par des modules partagés pour le chunking, la recherche,
l'évaluation et l'interface.

\begin{table}[H]
\centering
\caption{Composants logiciels principaux.}
\label{tab:outils}
\begin{tabularx}{\textwidth}{p{3.5cm}Y Y}
\toprule
\textbf{Composant} & \textbf{Technologie} & \textbf{Rôle}\\
\midrule
Langage & Python & Implémentation de l'ensemble du pipeline.\\
Index vectoriel & FAISS & Recherche sémantique des chunks.\\
Embeddings & multilingual-e5-base & Encodage des questions françaises et des passages.\\
Recherche lexicale & SimpleBM25 & Correspondance des termes techniques et APIs.\\
Reranking & bge-reranker-v2-m3 & Réordonnancement précis des candidats.\\
Génération & Mistral-7B-Instruct-v0.3 & Génération locale de réponses en contexte.\\
Évaluation & Ragas 0.4.3 & Mesures de récupération et de génération.\\
Interface & Gradio & Démonstration interactive dans Kaggle.\\
\bottomrule
\end{tabularx}
\end{table}

\section{Collecte, nettoyage et découpage}

Les documents Python, Scikit-learn et LangChain sont collectés depuis leurs
sources officielles~\cite{pythonDocs,sklearnDocs,langchainDocs}. Le nettoyage
supprime les éléments non utiles et conserve les informations nécessaires à la
provenance. Les documents nettoyés sont ensuite divisés en chunks de taille
contrôlée. Un chunk comprend son texte et des métadonnées telles que
\texttt{chunk\_id}, \texttt{document\_id}, \texttt{doc\_url} et une empreinte
de contenu.

L'identité stable des chunks est essentielle au protocole final. Elle rend
possible l'annotation des passages de référence et permet de distinguer une
mauvaise réponse causée par la génération d'une mauvaise réponse causée par la
récupération.

\section{Indexation et recherche hybride}

Lors de l'indexation, les chunks sont transformés en vecteurs. Les préfixes E5
sont respectés : \texttt{passage:} pour les chunks et \texttt{query:} pour les
questions. Les vecteurs sont normalisés avant la recherche de similarité.

La fonction de récupération interroge d'abord FAISS et BM25, fusionne les scores
selon un coefficient $\alpha=0{,}7$, récupère vingt candidats et applique le
reranker BGE. Les cinq premiers passages deviennent le contexte exact transmis
au modèle. Ce contexte est sauvegardé pendant l'évaluation afin de rendre les
scores auditables.

\begin{lstlisting}[caption={Principe simplifié de recherche hybride et reranking.},label={lst:retrieval}]
# La question est encodee avec le prefixe E5 approprie.
query_embedding = embedding_model.encode(["query: " + question])

# FAISS et BM25 produisent des scores complementaires.
hybrid_scores = alpha * dense_scores + (1 - alpha) * bm25_scores
candidates = top_k(hybrid_scores, k=20)

# Le cross-encoder reordonne seulement les 20 candidats.
final_contexts = reranker.rank(question, candidates)[:5]
\end{lstlisting}

\section{Génération de la réponse}

Le modèle Mistral est chargé à la demande par le client Hugging Face. Le prompt
contient la question, les passages numérotés et des règles explicites : répondre
uniquement à partir des passages, citer les identifiants de passages et ne pas
inventer d'API ou de paramètres absents du contexte. Une température faible est
utilisée afin de privilégier des réponses factuelles et stables.

\section{Interface Gradio}

L'interface permet de saisir une question, d'afficher la réponse et de consulter
les sources récupérées. Gradio est adapté à cette fonction de démonstration, car
la bibliothèque facilite la construction et le partage rapide d'interfaces de
modèles d'apprentissage automatique~\cite{abid2019gradio,gradio}. La même instance
du pipeline est réutilisée afin d'éviter le chargement multiple du modèle sur le
GPU. Lorsque le pipeline détecte un sujet explicitement hors périmètre,
l'interface affiche un statut clair au lieu de montrer des sources potentiellement
trompeuses.

\begin{figure}[H]
\centering
\begin{tikzpicture}[node distance=1.2cm and 1.1cm,>=Latex,
  box/.style={draw,rounded corners=4pt,fill=green!8,minimum width=3.3cm,
  minimum height=0.9cm,align=center,font=\small}]
\node[box] (q) {Question utilisateur};
\node[box,right=of q] (check) {Contrôle de périmètre};
\node[box,right=of check] (r) {Réponse sourcée};
\node[below=0.9cm of check,draw,rounded corners=4pt,fill=red!7,minimum width=3.5cm,
 align=center,font=\small] (refuse) {Refus clair : corpus insuffisant\\ou bibliothèque non couverte};
\draw[->,thick] (q)--(check);
\draw[->,thick] (check)--node[above,font=\scriptsize]{contexte suffisant}(r);
\draw[->,thick] (check)--node[right,font=\scriptsize]{hors périmètre}(refuse);
\end{tikzpicture}
\caption{Comportement de l'interface face à une question hors périmètre.}
\label{fig:guard}
\end{figure}

\section{Protection contre les questions hors périmètre}

Le cas \texttt{pd.head()} a montré qu'un système spécialisé ne doit pas répondre
à une bibliothèque non incluse dans son corpus. Une réponse antérieure avait
produit une liste de noms hors sujet, car le générateur avait reçu des passages
non pertinents. La correction ajoutée vérifie donc des alias et des noms de
bibliothèques explicitement non couverts avant toute génération.

\begin{lstlisting}[caption={Logique simplifiée de protection hors périmètre.},label={lst:guard}]
def explicit_scope_refusal(question, config):
    topic = detect_unsupported_library(question, config)
    if topic:
        return {
            "allow_answer": False,
            "status": "hors_perimetre",
            "message": "Bibliotheque absente du corpus actuel."
        }
    return None

refusal = explicit_scope_refusal(question, scope_config)
if refusal is not None:
    return refusal              # Mistral n'est pas appele
contexts = retrieve(question)   # seulement pour une question couverte
\end{lstlisting}

Le seuil numérique de score est conservé comme amélioration future. Les scores
de reranking ne sont pas des probabilités universelles ; ils doivent être
calibrés avec un jeu dédié de questions dans et hors périmètre avant de servir
à un refus automatique. Cette décision évite de rejeter arbitrairement des
questions valides en français portant sur une documentation anglaise.

\section{Notebook Kaggle et sauvegarde des artefacts}

Le notebook français pilote les étapes du pipeline et distingue deux modes :
\texttt{PREPARATION} pour créer les artefacts et \texttt{EVALUATION\_FINALE}
pour exécuter Ragas sur le jeu de 120 questions. Les résultats sont exportés
sous forme de JSON, CSV, graphiques et archive ZIP. Les clés API ne sont pas
insérées dans le notebook : elles sont lues depuis les secrets Kaggle.

\section*{Conclusion du chapitre}
\addcontentsline{toc}{section}{Conclusion du Chapitre 4}

La réalisation met en œuvre l'ensemble du cycle RAG et ajoute des mécanismes de
traçabilité et de protection. Le chapitre suivant présente l'évaluation des
réponses obtenues, en distinguant les résultats valides des limites
méthodologiques encore présentes.

%=============================================================================
\chapter{Évaluation expérimentale et analyse des résultats}
\label{chap:evaluation}

\section{Protocole expérimental}

L'évaluation comprend deux niveaux complémentaires. Le premier est une baseline
historique, archivée séparément afin de conserver le diagnostic du système avant
les améliorations. Le second est le protocole principal, construit pour une
évaluation Ragas plus large, traçable et équilibrée.

\subsection{Baseline historique : rôle et limites de comparaison}

La baseline \texttt{baseline\_v1\_custom\_ragas\_inspired} a été exécutée sur
20 questions : 7 portant sur Python, 6 sur Scikit-learn et 7 sur LangChain. Son
artefact est conservé au commit \texttt{c88e287}~\cite{baselineArtifact,baselineCode}.
Elle utilise un évaluateur personnalisé inspiré de Ragas, avec au plus cinq
contextes et huit affirmations examinées par réponse. Elle constitue donc un
\textbf{diagnostic initial}, et non l'évaluation officielle finale.

Le code historique prévoit comme configuration de repli un récupérateur
\texttt{all-MiniLM-L6-v2}, issu de la famille compacte MiniLM~\cite{wang2020minilm},
une recherche sémantique seule et cinq passages.
Cependant, il pouvait charger un rapport de benchmarking présent dans Kaggle ;
ce manifeste précis n'a pas été archivé avec l'exécution historique. Par rigueur,
le modèle MiniLM est présenté comme une configuration de repli du code et non
comme un paramètre expérimental certifié du run baseline~\cite{baselineCode}.

La baseline est utile car elle a révélé, sur un jeu réduit, les premières limites
de pertinence et de récupération. Elle a orienté les améliorations ultérieures :
provenance des chunks, jeu de test équilibré, recherche E5--BM25, reranking et
évaluation Ragas officielle. Elle ne peut pas servir à calculer un gain chiffré
direct, car les questions, les références de contexte, l'implémentation des
métriques et le juge diffèrent du protocole final.

\subsection{Protocole final}

L'évaluation principale utilise Ragas 0.4.3 sur 120 questions françaises,
équilibrées entre Python, Scikit-learn et LangChain. Le jeu final contient une
réponse de référence, des URLs de documentation et des identifiants de chunks
pour chaque question. La validation des références est tracée comme une
validation experte assistée par IA à relecteur unique. Cette précision doit être
déclarée : elle ne constitue pas une annotation indépendante par plusieurs
évaluateurs humains.

\begin{table}[H]
\centering
\caption{Comparaison des protocoles de la baseline et de l'évaluation finale.}
\label{tab:protocoles}
\begin{tabularx}{\textwidth}{p{3.5cm}Y Y}
\toprule
\textbf{Caractéristique} & \textbf{Baseline historique} & \textbf{Évaluation finale}\\
\midrule
Nombre de questions & 20 & 120 (40 par domaine).\\
Évaluateur & Code personnalisé inspiré de Ragas. & Ragas 0.4.3 officiel avec juge Mistral.\\
Modèle / recherche & Configuration historique non entièrement gelée ; repli du code : all-MiniLM-L6-v2, sémantique, top 5. & E5 multilingue + BM25, $\alpha=0{,}7$, reranking BGE, 20 candidats puis top 5.\\
Références de contexte & Non systématiquement annotées. & IDs de chunks et URLs de référence.\\
Rôle & Diagnostic initial. & Évaluation principale du système final.\\
Comparaison numérique & Non comparable strictement. & Protocole plus large et plus traçable.\\
\bottomrule
\end{tabularx}
\end{table}

\section{Configuration de l'évaluation finale}

Le générateur local est \texttt{mistralai/Mistral-7B-Instruct-v0.3}. Le juge
Ragas est \texttt{mistral-small-latest}, appelé via l'interface asynchrone
compatible OpenAI de Mistral. Cette séparation évite de confondre le modèle qui
répond à l'utilisateur et le modèle qui juge la réponse.

Le run analysé est identifié par \texttt{final\_120\_20260827T154025Z}. Le
runner sauvegarde la question, la réponse, les chunks récupérés, le contexte
exact transmis au générateur, les scores, les raisons et les erreurs éventuelles.
Cette granularité permet un audit par question, au-delà des seules moyennes.

\section{Couverture de l'évaluation}

Le passage au juge Mistral a permis de calculer les métriques Ragas pour la
quasi-totalité des échantillons. Sur 120 questions, 114 sont entièrement
évaluées et 6 sont partielles. Aucun échantillon n'est en échec complet.

\begin{table}[H]
\centering
\caption{État d'exécution des 120 évaluations finales.}
\label{tab:couverture}
\begin{tabular}{lrr}
\toprule
\textbf{État} & \textbf{Nombre} & \textbf{Part}\\
\midrule
Échantillons entièrement évalués & 114 & 95,0\%\\
Échantillons partiellement évalués & 6 & 5,0\%\\
Échecs complets & 0 & 0,0\%\\
\bottomrule
\end{tabular}
\end{table}

Les six valeurs manquantes résultent de réponses incomplètes du juge (limite de
tokens) ou d'indisponibilités temporaires du service Mistral. Elles ne doivent
pas être transformées artificiellement en score nul. Le runner conserve la
valeur absente et enregistre l'erreur, ce qui préserve l'intégrité du résultat.

\section{Résultats Ragas globaux}

Les résultats globaux de l'exécution sont présentés au tableau~\ref{tab:ragasglobal}.
Un score proche de 1 est favorable. Les taux de couverture indiquent la part des
questions pour lesquelles une valeur a effectivement été calculée.

\begin{table}[H]
\centering
\caption{Résultats Ragas globaux sur 120 questions.}
\label{tab:ragasglobal}
\begin{tabularx}{\textwidth}{p{3.2cm}p{2.2cm}p{2.6cm}Y}
\toprule
\textbf{Métrique} & \textbf{Moyenne} & \textbf{Couverture} & \textbf{Interprétation}\\
\midrule
Faithfulness & \textbf{0,888} & 116/120 & Les réponses sont globalement justifiées par les passages fournis.\\
Answer Relevancy & \textbf{0,913} & 120/120 & Les réponses traitent le plus souvent la question posée.\\
Context Precision & \textbf{0,819} & 119/120 & Les passages sont souvent utiles, mais le classement peut être amélioré.\\
Context Recall & \textbf{0,879} & 120/120 & Les informations nécessaires sont fréquemment présentes dans le contexte.\\
Factual\newline Correctness & \textbf{0,693} & 118/120 & L'exactitude par rapport aux références est l'axe principal d'amélioration.\\
\bottomrule
\end{tabularx}
\end{table}

Le système présente un bon niveau d'ancrage documentaire : Faithfulness et
Answer Relevancy sont élevées. Cependant, Factual Correctness est sensiblement
plus basse. Ce résultat signifie qu'une réponse peut être pertinente et
soutenue par un contexte, tout en restant imprécise, trop générale ou incomplète
par rapport à la réponse de référence.

\section{Résultats par domaine}

\begin{table}[H]
\centering
\caption{Moyennes Ragas par domaine.}
\label{tab:ragasdomaines}
\begin{tabular}{lrrrrr}
\toprule
\textbf{Domaine} & \textbf{Faith.} & \textbf{Rel.} & \textbf{C. Prec.} & \textbf{C. Rec.} & \textbf{Fact. Corr.}\\
\midrule
Python & 0,890 & 0,910 & \textbf{0,869} & 0,875 & \textbf{0,758}\\
Scikit-learn & \textbf{0,891} & \textbf{0,924} & 0,864 & \textbf{0,925} & 0,708\\
LangChain & 0,882 & 0,907 & 0,726 & 0,838 & 0,611\\
\bottomrule
\end{tabular}
\end{table}

Python obtient la meilleure exactitude factuelle, tandis que Scikit-learn
présente le meilleur rappel de contexte. LangChain est le domaine le plus
difficile : la précision des contextes, le rappel et l'exactitude factuelle y
sont les plus faibles. Cette observation est cohérente avec une documentation
plus hétérogène, comportant de nombreuses pages d'intégration et des concepts
proches lexicalement mais distincts fonctionnellement.

\section{Diagnostic déterministe de récupération}

Les identifiants de chunks annotés permettent de calculer des métriques de
récupération indépendantes du juge LLM. Un Hit@5 de 0,792 indique qu'au moins
un chunk de référence est présent parmi les cinq passages récupérés pour 95
questions sur 120. Le MRR@5 de 0,633 indique que le premier chunk de référence
est souvent bien classé, mais pas systématiquement premier.

\begin{table}[H]
\centering
\caption{Diagnostic déterministe de récupération.}
\label{tab:retrievaldiag}
\begin{tabular}{lrrr}
\toprule
\textbf{Domaine} & \textbf{Hit@5} & \textbf{MRR@5} & \textbf{Référence absente du top 5}\\
\midrule
Python & 0,800 & 0,618 & 8/40\\
Scikit-learn & \textbf{0,925} & \textbf{0,754} & 3/40\\
LangChain & \textbf{0,650} & \textbf{0,527} & 14/40\\
\midrule
Global & 0,792 & 0,633 & 25/120\\
\bottomrule
\end{tabular}
\end{table}

La précision déterministe par identifiants est faible (environ 0,240), mais elle
doit être interprétée avec précaution. Le système retourne systématiquement cinq
chunks, alors que de nombreuses questions n'ont qu'un seul chunk de référence.
Même avec ce chunk parfaitement classé en première position, la précision par
identifiants est alors mécaniquement égale à $1/5=0{,}20$. Hit@5, MRR@5 et
rappel constituent donc des compléments indispensables.

\section{Analyse qualitative de cas faibles}

L'analyse des échantillons individuels confirme que les moyennes cachent des
comportements variés. La question sur la différence entre LangChain et LangGraph
présente une faible pertinence et une faible précision de contexte. La question
sur la compression automatique de contexte dans Deep Agents présente une
fidélité nulle. D'autres cas faibles concernent les annotations de types Python,
le choix de scaler Scikit-learn en présence de valeurs aberrantes et le choix de
métrique avant l'optimisation d'hyperparamètres.

Ces observations conduisent à une conclusion opérationnelle : les améliorations
futures doivent prioriser le corpus et le classement LangChain, plutôt que de
remplacer immédiatement le générateur. La qualité de la réponse dépend d'abord
de la présence d'un passage pertinent dans le top 5.

\section{Limite identifiée dans le jeu de test}

Un contrôle post-exécution a détecté cinq lignes mal formées dans le jeu de test
final. Dans ces enregistrements, le champ question contient une phrase
déclarative proche d'une réponse, tandis que le champ de référence contient une
URL au lieu d'une réponse textuelle. Les identifiants concernés sont
\texttt{test\_python\_020}, \texttt{test\_python\_022},
\texttt{test\_python\_024}, \texttt{test\_python\_032} et
\texttt{test\_scikit\_learn\_003}.

Cette anomalie concerne 4,2\% du jeu. Une analyse de sensibilité effectuée sur
les 115 lignes correctement formées donne une vision plus prudente des résultats
provisoires.

\begin{table}[H]
\centering
\caption{Analyse de sensibilité sur les 115 questions correctement formées.}
\label{tab:sensibilite}
\begin{tabular}{lr}
\toprule
\textbf{Métrique} & \textbf{Moyenne sur 115 lignes}\\
\midrule
Faithfulness & 0,888\\
Answer Relevancy & 0,916\\
Context Precision & 0,847\\
Context Recall & 0,917\\
Factual Correctness & 0,698\\
\bottomrule
\end{tabular}
\end{table}

Les conclusions principales restent stables, mais une présentation académique
rigoureuse doit signaler cette limite. Les cinq lignes doivent être corrigées et
le run doit être reproduit une fois avant de qualifier les valeurs de scores
finals définitivement gelés.

\section{Positionnement par rapport à la baseline}

L'artefact de baseline enregistre les moyennes suivantes : Faithfulness = 0,789,
Answer Relevancy = 0,727, Context Precision = 0,562 et Context Recall = 0,625.
Il contient aussi des indicateurs de récupération élevés sur son jeu court :
Hit rate@5 = 0,980, MRR@5 = 0,927 et Precision@5 = 0,825~\cite{baselineArtifact}.
Ces résultats ont eu une utilité exploratoire : ils ont aidé à formuler les
axes de travail et à décider de renforcer la traçabilité, la diversité du jeu de
questions et l'évaluation des contextes.

Ils ne doivent cependant pas être comparés directement aux scores Ragas finaux.
La baseline repose sur 20 questions, un évaluateur personnalisé, des références
de contexte incomplètement annotées et une configuration historique dont le
manifeste de benchmarking n'est plus disponible. L'évaluation finale porte sur
120 questions équilibrées, des réponses de référence, des URLs, des chunks
annotés, Ragas 0.4.3 et un juge Mistral. Les deux expériences ne mesurent donc
pas exactement le même objet~\cite{baselineArtifact,es2023ragas}.

La conclusion méthodologique est la suivante : la baseline a servi de
\textbf{diagnostic initial}, tandis que l'évaluation finale constitue le résultat
principal du projet. La progression démontrée ici est une progression de la
rigueur expérimentale -- corpus traçable, récupération hybride, reranking et
protocole documenté -- et non une affirmation de gain numérique strict entre
deux protocoles non identiques.

\section*{Conclusion du chapitre}
\addcontentsline{toc}{section}{Conclusion du Chapitre 5}

L'évaluation finale fournit des résultats Ragas majoritairement calculés et
traçables. Les scores élevés de fidélité et de pertinence constituent des signes
encourageants. Les faiblesses observées portent principalement sur l'exactitude
factuelle et sur la récupération dans le domaine LangChain. Les limites du jeu
de test et du juge sont explicitement documentées, ce qui renforce la valeur
scientifique de l'analyse.

%=============================================================================
\chapter{Discussion, limitations et perspectives}
\label{chap:discussion}

\section{Bilan des contributions}

Le projet a permis de construire un pipeline RAG complet et reproductible pour
un corpus documentaire technique. Les contributions principales sont la
traçabilité stable des chunks, l'association entre recherche dense et lexicale,
le reranking multilingue, l'interface Gradio, le jeu d'évaluation de 120
questions et la sauvegarde d'artefacts détaillés d'exécution.

La valeur du travail ne réside pas uniquement dans les scores moyens, mais dans
la capacité à expliquer leur origine. Le manifeste d'index, le contexte exact du
prompt, les identifiants de chunks et les erreurs par question rendent possible
une analyse qui sépare récupération, génération et disponibilité des services
externes.

\section{Difficultés rencontrées et résolution}

Plusieurs difficultés techniques ont été rencontrées lors de l'évaluation. La
première concernait des incompatibilités d'interface entre Ragas et les classes
d'embeddings : certaines méthodes asynchrones étaient exigées par Ragas. Une
implémentation concrète de l'interface moderne d'embeddings a été ajoutée et
couverte par des tests.

La seconde difficulté concernait le juge OpenAI, dont le quota API était nul.
Le juge a alors été remplacé par Mistral via son interface compatible OpenAI et
un client asynchrone. Cette adaptation a permis de calculer les métriques Ragas
pour l'essentiel des 120 questions. Enfin, des erreurs temporaires de réseau ou
de sortie incomplète du juge ont été conservées de façon transparente plutôt que
remplacées par des valeurs arbitraires.

\section{Limitations}

\begin{table}[H]
\centering
\caption{Limitations et réponses proposées.}
\label{tab:limitations}
\begin{tabularx}{\textwidth}{p{4.0cm}Y Y}
\toprule
\textbf{Limitation} & \textbf{Impact} & \textbf{Perspective}\\
\midrule
Cinq lignes de test mal formées & Les scores à 120 lignes restent provisoires. & Corriger les enregistrements puis relancer une seule évaluation finale.\\
Validation experte par IA unique & Variabilité possible de l'annotation des références. & Ajouter une revue humaine ou un second évaluateur indépendant.\\
Faiblesse LangChain & 14 références absentes du top 5. & Enrichir les sources, filtrer par section et ajuster le reranking.\\
Questions hors périmètre & Risque de réponse trompeuse. & Garde déterministe et calibration ultérieure d'un seuil de confiance.\\
Juge API externe & Risque de quota ou d'indisponibilité. & Prévoir un fournisseur de secours et conserver les artefacts d'erreurs.\\
\bottomrule
\end{tabularx}
\end{table}

\section{Perspectives d'amélioration}

La première priorité est la correction des cinq lignes du jeu de test, suivie
d'une unique réexécution finale de Ragas. La deuxième priorité est
l'amélioration ciblée de LangChain. Elle peut reposer sur l'ajout de pages
spécifiques, la normalisation des URLs, des filtres par domaine ou par section,
et une étude des questions dont la référence est absente du top 5.

La troisième priorité concerne la robustesse de l'interface. La garde hors
périmètre peut être étendue par un seuil calibré sur un jeu de questions dans et
hors corpus. L'interface pourrait aussi afficher les extraits de passages et
une estimation de confiance, afin que l'utilisateur comprenne le fondement de
la réponse.

Enfin, une évaluation plus robuste pourrait associer le juge automatique à une
revue humaine d'un sous-échantillon. La comparaison entre plusieurs juges et
plusieurs annotateurs permettrait d'étudier la stabilité des métriques et de
renforcer la validité externe des conclusions.

%=============================================================================
\chapter*{Conclusion générale}
\addcontentsline{toc}{chapter}{Conclusion générale}
\markboth{Conclusion générale}{}

Ce Projet de Fin d'Année a porté sur la conception et l'évaluation d'un assistant
RAG dédié à la documentation technique de Python, Scikit-learn et LangChain. Le
système combine une collecte de sources officielles, un nettoyage et un
chunking traçables, une indexation FAISS, une recherche hybride BM25--E5, un
reranking BGE et une génération par Mistral.

La mise en place d'un protocole à 120 questions, de références documentaires et
d'identifiants de chunks constitue une contribution méthodologique importante.
L'exécution Ragas avec un juge Mistral a produit une fidélité au contexte de
0,888 et une pertinence de réponse de 0,913, ce qui indique que les réponses
sont généralement reliées aux passages récupérés et aux questions formulées.
L'exactitude factuelle de 0,693 révèle toutefois qu'une réponse pertinente ne
suffit pas : la précision par rapport à une référence reste un enjeu central.

Le travail met également en évidence des limites clairement identifiées :
LangChain est moins bien récupéré que les deux autres domaines, cinq lignes du
jeu final doivent être corrigées et une validation à plusieurs annotateurs
renforcerait la démarche. Ces limites ne diminuent pas l'intérêt du projet ;
elles constituent au contraire une base précise pour une amélioration future,
plus contrôlée et plus mesurable.

En conclusion, le projet montre qu'un assistant RAG spécialisé peut offrir un
accès plus structuré et plus sourcé à une documentation technique. Sa valeur
professionnelle repose sur la combinaison d'une architecture opérationnelle,
d'une interface de démonstration et d'une évaluation reproductible qui ne
masque ni les performances observées ni les limites restantes.

%=============================================================================
% BIBLIOGRAPHIE IEEE INTÉGRÉE
%=============================================================================
\clearpage
\addcontentsline{toc}{chapter}{Bibliographie}
\begin{thebibliography}{99}

\bibitem{lewis2020rag}
P. Lewis, E. Perez, A. Piktus et al., ``Retrieval-Augmented Generation for
Knowledge-Intensive NLP Tasks,'' \textit{Advances in Neural Information
Processing Systems}, vol. 33, 2020. Disponible : \url{https://arxiv.org/abs/2005.11401}.

\bibitem{reimers2019sbert}
N. Reimers et I. Gurevych, ``Sentence-BERT: Sentence Embeddings using Siamese
BERT-Networks,'' \textit{Proceedings of EMNLP-IJCNLP}, 2019. Disponible :
\url{https://arxiv.org/abs/1908.10084}.

\bibitem{johnson2017faiss}
J. Johnson, M. Douze et H. Jégou, ``Billion-scale similarity search with
GPUs,'' \textit{arXiv preprint arXiv:1702.08734}, 2017. Disponible :
\url{https://arxiv.org/abs/1702.08734}.

\bibitem{robertson2009bm25}
S. Robertson et H. Zaragoza, ``The Probabilistic Relevance Framework: BM25 and
Beyond,'' \textit{Foundations and Trends in Information Retrieval}, vol. 3,
no. 4, pp. 333--389, 2009.

\bibitem{nogueira2019passagererank}
R. Nogueira et K. Cho, ``Passage Re-ranking with BERT,'' \textit{arXiv preprint
arXiv:1901.04085}, 2019. Disponible : \url{https://arxiv.org/abs/1901.04085}.

\bibitem{wange5}
L. Wang, N. Yang, X. Huang et al., ``Text Embeddings by Weakly-Supervised
Contrastive Pre-training,'' \textit{arXiv preprint arXiv:2212.03533}, 2022.
Disponible : \url{https://arxiv.org/abs/2212.03533}.

\bibitem{jiang2023mistral}
A. Q. Jiang, A. Sablayrolles, A. Roux et al., ``Mistral 7B,'' \textit{arXiv
preprint arXiv:2310.06825}, 2023. Disponible : \url{https://arxiv.org/abs/2310.06825}.

\bibitem{es2023ragas}
S. Es, J. James, L. Espinosa-Anke et S. Schockaert, ``Ragas: Automated
Evaluation of Retrieval Augmented Generation,'' \textit{arXiv preprint
arXiv:2309.15217}, 2023. Disponible : \url{https://arxiv.org/abs/2309.15217}.

\bibitem{ragasdocs}
Ragas, ``Available Metrics,'' documentation officielle, consultée en août 2026.
Disponible : \url{https://docs.ragas.io/en/stable/concepts/metrics/available_metrics/}.

\bibitem{pythonDocs}
Python Software Foundation, ``Python 3 Documentation,'' documentation officielle,
consultée en août 2026. Disponible : \url{https://docs.python.org/3/}.

\bibitem{sklearnDocs}
Scikit-learn Developers, ``Scikit-learn User Guide,'' documentation officielle,
consultée en août 2026. Disponible : \url{https://scikit-learn.org/stable/user_guide.html}.

\bibitem{langchainDocs}
LangChain, ``LangChain Documentation,'' documentation officielle, consultée en
 août 2026. Disponible : \url{https://python.langchain.com/docs/}.

\bibitem{bgereranker}
Beijing Academy of Artificial Intelligence, ``bge-reranker-v2-m3,'' fiche modèle,
consultée en août 2026. Disponible : \url{https://huggingface.co/BAAI/bge-reranker-v2-m3}.

\bibitem{mistraldocs}
Mistral AI, ``Python SDK and API documentation,'' documentation officielle,
consultée en août 2026. Disponible : \url{https://docs.mistral.ai/resources/sdks}.

\bibitem{gradio}
Gradio, ``Gradio Documentation,'' documentation officielle, consultée en août
2026. Disponible : \url{https://www.gradio.app/docs}.

\bibitem{kaggle}
Kaggle, ``Notebooks Documentation,'' documentation officielle, consultée en août
2026. Disponible : \url{https://www.kaggle.com/docs/notebooks}.

\bibitem{karpukhin2020dpr}
V. Karpukhin, B. Oğuz, S. Min et al., ``Dense Passage Retrieval for Open-Domain
Question Answering,'' \textit{Proceedings of EMNLP}, 2020. Disponible :
\url{https://arxiv.org/abs/2004.04906}.

\bibitem{gao2023survey}
Y. Gao, Y. Xiong, X. Gao et al., ``Retrieval-Augmented Generation for Large
Language Models: A Survey,'' \textit{arXiv preprint arXiv:2312.10997}, 2023.
Disponible : \url{https://arxiv.org/abs/2312.10997}.

\bibitem{yu2024eval}
H. Yu, A. Gan, K. Zhang et al., ``Evaluation of Retrieval-Augmented Generation:
A Survey,'' \textit{arXiv preprint arXiv:2405.07437}, 2024. Disponible :
\url{https://arxiv.org/abs/2405.07437}.

\bibitem{vaswani2017attention}
A. Vaswani, N. Shazeer, N. Parmar et al., ``Attention Is All You Need,''
\textit{Advances in Neural Information Processing Systems}, vol. 30, 2017.
Disponible : \url{https://arxiv.org/abs/1706.03762}.

\bibitem{devlin2018bert}
J. Devlin, M.-W. Chang, K. Lee et K. Toutanova, ``BERT: Pre-training of Deep
Bidirectional Transformers for Language Understanding,'' \textit{Proceedings of
NAACL-HLT}, 2019. Disponible : \url{https://arxiv.org/abs/1810.04805}.

\bibitem{wolf2019transformers}
T. Wolf, L. Debut, V. Sanh et al., ``Transformers: State-of-the-Art Natural
Language Processing,'' \textit{Proceedings of EMNLP}, 2020. Disponible :
\url{https://arxiv.org/abs/1910.03771}.

\bibitem{abid2019gradio}
A. Abid, A. Abdalla, A. Abid et al., ``Gradio: Hassle-Free Sharing and Testing
of ML Models in the Wild,'' \textit{ICML Workshop on Human in the Loop Learning},
2019. Disponible : \url{https://arxiv.org/abs/1906.02569}.

\bibitem{wang2020minilm}
W. Wang, F. Wei, L. Dong et al., ``MiniLM: Deep Self-Attention Distillation for
Task-Agnostic Compression of Pre-Trained Transformers,'' \textit{arXiv preprint
arXiv:2002.10957}, 2020. Disponible : \url{https://arxiv.org/abs/2002.10957}.

\bibitem{baselineArtifact}
E. Dahani, ``Baseline RAG historique : métriques archivées,'' artefact interne
du projet, fichier \texttt{evaluation/baseline/baseline\_metrics\_v1.json},
commit source \texttt{c88e287}, août 2026.

\bibitem{baselineCode}
E. Dahani, ``Évaluateur baseline personnalisé inspiré de Ragas,'' code source
archivé, fichier \texttt{legacy/etape6\_baseline\_custom.py}, commit source
\texttt{c88e287}, août 2026.

\bibitem{organisationSource}
3D Smart Factory, ``Présentation de l'organisme d'accueil,'' informations
internes communiquées dans les deux extraits de rapport remis pour le stage,
Mohammédia, août 2026.

\end{thebibliography}

%=============================================================================
% ANNEXES
%=============================================================================
\appendix
\chapter{Reproductibilité et artefacts sauvegardés}

Le projet est versionné dans un dépôt GitHub. Le notebook Kaggle français
sépare le code, les données et les résultats afin d'éviter qu'une mise à jour du
code efface le corpus ou l'index. Les artefacts de l'évaluation finale comprennent
notamment \texttt{samples.jsonl}, \texttt{samples.csv}, \texttt{summary.json},
\texttt{manifest.json}, les graphiques de métriques et une archive
\texttt{resultats\_rag\_stage.zip}.

\begin{longtable}{p{4.2cm}p{10cm}}
\toprule
\textbf{Artefact} & \textbf{Utilité}\\
\midrule
\endhead
\texttt{index\_manifest.json} & Configuration de l'index, profil de récupération et modèle d'embeddings.\\
\texttt{test\_dataset\_v1.\allowbreak jsonl} & Questions, références, URLs et IDs de chunks du jeu final.\\
\texttt{samples.jsonl} & Données par question : contextes, réponse, scores, raisons et erreurs.\\
\texttt{summary.json} & Moyennes, couverture des métriques et résultats par domaine.\\
\texttt{manifest.json} & Informations de reproductibilité sur le run et le juge.\\
\texttt{resultats\_rag\_stage.\allowbreak zip} & Archive exportable regroupant les principaux résultats.\\
\bottomrule
\end{longtable}

\chapter{Schéma conceptuel d'un enregistrement de test}

Chaque question du jeu final suit une structure JSONL similaire à celle-ci :

\begin{lstlisting}[caption={Structure simplifiée d'un enregistrement du jeu de test.}]
{
  "question_id": "test_python_001",
  "split": "test",
  "domain": "python",
  "language": "fr",
  "user_input": "Quelle est la portee d'une variable ?",
  "reference": "La portee indique ou une variable est accessible.",
  "reference_context_ids": ["python_..._chunk_0001"],
  "reference_source_urls": ["https://docs.python.org/3/..."],
  "annotation": {"review_status": "validated"}
}
\end{lstlisting}

\chapter{Éléments à compléter avant dépôt officiel}

Avant la soumission officielle, il est recommandé de vérifier les éléments
suivants : les noms et fonctions exacts des encadrants sur la page de garde, la
date de soutenance, le numéro éventuel du PFA, l'orthographe des membres du jury
et les règles formelles spécifiques de l'INSEA. Les cinq lignes mal formées du
jeu de test doivent aussi être corrigées avant de présenter les métriques à 120
questions comme des valeurs définitivement gelées.

\end{document}
